\documentclass[12pt]{article}
\usepackage{geometry}
\geometry{letterpaper, margin=1in}
\usepackage{hyperref}
\usepackage[usenames,dvipsnames]{xcolor}
\usepackage{microtype}
\usepackage{enumitem}

\definecolor{accent}{RGB}{26, 60, 110}

\begin{document}

%----------HEADER----------
\begin{center}
  {\LARGE\textbf{\textcolor{accent}{SITONG LI}}} \\[4pt]
  \small Minneapolis, MN
  \;\textbullet\;
  \href{mailto:ruc.sitongli@gmail.com}{ruc.sitongli@gmail.com}
  \;\textbullet\;
  (952) 201-4026
\end{center}

\vspace{1em}

\noindent May 15, 2026

\vspace{1em}

\noindent Dear Prof.\ Kanter and the USC Schaeffer Center Hiring Committee,

\vspace{1em}

I am writing to apply for the Project Specialist position supporting Professor Genevieve Kanter's research at the USC Leonard D.\ Schaeffer Center for Health Policy and Economics. I am currently completing a full-time predoctoral fellowship at the Ross School of Business, University of Michigan, working under faculty mentors with appointments at Michigan, Yale, and Barnard on projects analyzing large-scale administrative data. That experience has confirmed what my prior research assistant work suggested: rigorous academic economics research is the career I am fully committed to, and I am actively preparing for a Ph.D.\ in economics. The Schaeffer Center's mission --- producing evidence that shapes health policy decisions --- is exactly the kind of policy-relevant applied economics I want to spend my career on, and I am eager to bring the research skills I have built to Professor Kanter's program.

At Michigan Ross, I work full-time as a predoctoral research fellow under Francesca Truffa (Michigan), Menaka Hampole (Yale), and Ashley Wong (Barnard). My work centers on two large-scale data infrastructure projects: building a scalable entity resolution pipeline --- integrating LLMs, fuzzy string matching, and word embeddings --- to accurately link student application records across heterogeneous administrative databases, and designing a computer vision pipeline to infer demographic attributes from applicant photos for structured population analysis. Working full-time in a top research environment, embedded in a faculty-driven project at the publication stage, has sharpened my ability to independently manage ambiguous data problems, maintain rigorous analytical documentation, and deliver reliable outputs under academic production timelines. It has also deepened my commitment to the academic path: this is the kind of work I want to do, and I want to do it at the doctoral level.

My most directly relevant methodological experience for the Schaeffer Center's work is from a research assistantship with Prof.\ Haiwen Zhang at the University of Minnesota on a project analyzing EPA regulatory enforcement. I designed a data pipeline to convert ten years of scanned regulatory correspondence into structured, analysis-ready text using OCR and regular expressions, then applied an agentic GPT-4 pipeline to extract policy topics and inquiry points from over 100,000 emails across more than 1,000 regulated entities. I built a search bot integrated with the Google Search API to retrieve entity-specific contextual information, matched the email data with EPA enforcement records, and ran probit regressions to examine the causal determinants of enforcement actions. This project --- systematic online data collection, NLP-based document analysis, and causal econometric methods in service of a regulatory policy question --- is a close methodological analogue to health policy research of the kind Professor Kanter conducts.

My other research experience adds complementary depth. In a concurrent project with Prof.\ Zhang, I collected 10,000+ full-text SEC 10-K filings via web crawler, segmented documents using regex and HTML parsing, and applied BERT classifiers to analyze 500,000+ sentences, validating findings through panel regressions and event study models in Stata. As a research assistant to Prof.\ William Cassidy at Washington University in St.\ Louis, I used agentic GPT-4 pipelines to extract fiscal policy indicators from 100,000+ lines of congressional testimony; the resulting analyses contribute to a working paper now in preparation. Across all three settings --- plus my current predoctoral fellowship --- I have taken ownership of the full research production cycle: designing collection protocols, writing extraction pipelines, cleaning and merging heterogeneous datasets, running statistical models, and preparing analytical outputs for both academic and practitioner audiences.

I am proficient in Python, R, and Stata, and have hands-on deployment experience with agentic AI systems (GPT-4, BERT, Donut) applied to real research tasks. I am comfortable with academic database searching, citation and annotation management, and preparing summary reports and visualizations. I am available to begin immediately and am open to any work arrangement. I would welcome the opportunity to discuss how my background fits Professor Kanter's current research agenda.

\vspace{1em}

\noindent Sincerely,\\[0.5em]
\textbf{Sitong Li}

\end{document}
